\hitFrontHeading{Abstract}

Efficient energy systems are moving from experience-driven operation toward data-driven planning, scheduling, and maintenance. In integrated energy stations, combined heat and power systems, and distributed storage devices, sensor measurements, operation logs, weather forecasts, and equipment records usually differ in sampling frequency, noise level, and semantic granularity. This public example builds a doctoral dissertation scaffold around multi-source data fusion, state modeling, scheduling optimization, and result interpretation, so that the `thesis-hit-doctor` template can demonstrate the complete structure required by Harbin Institute of Technology.

The example first organizes heterogeneous data through a layered quality-control framework and maps equipment states, environmental disturbances, and operation constraints into a computable feature space. It then constructs an energy-efficiency prediction model to describe system behavior under varying load conditions. A rolling scheduling method is further introduced to balance energy consumption, peak load, and operation stability while respecting device safety constraints. Finally, an interpretation workflow is provided so that the results can be reviewed in engineering practice and documented in a dissertation. All data and narratives are public demonstration materials and do not contain private or confidential information.

At the typesetting level, the template covers the Chinese cover, Chinese title page, English title page, Chinese and English abstracts, table of contents, English contents, chapters, figures, tables, equations, conclusions, references, innovative achievements, review and defense information, originality statement, acknowledgements, and resume. The whole project is rendered by the unified `bensz-thesis` build pipeline and can be refined against the official HIT Word example in later visual validation.

\hitKeywordsEnLine
